% Draft 7 (revised), 12 November 2026.
% Bibliography: references.bib, shared with the companion paper hodge_classes_algebraic_cycles.tex.

\documentclass[11pt,reqno]{amsart}

% ----------------------------------------------------------------------
% Packages
% ----------------------------------------------------------------------
\usepackage[T1]{fontenc}
\usepackage{lmodern}
\usepackage{amssymb}
\usepackage{microtype}
\usepackage[a4paper,margin=1.15in]{geometry}
\usepackage{enumitem}
\usepackage{xcolor}
\usepackage{hyperref}
\usepackage[nameinlink,noabbrev]{cleveref}

\hypersetup{
  colorlinks=true,
  linkcolor=blue!60!black,
  citecolor=blue!60!black,
  urlcolor=blue!60!black,
  pdftitle={Notes on Analytification, GAGA, Kaehler Identities, and Lefschetz Theory}
}

\setlist[enumerate]{label=(\roman*),leftmargin=2.2em}

% ----------------------------------------------------------------------
% Theorem environments
% ----------------------------------------------------------------------
\newtheorem{theorem}{Theorem}[section]
\newtheorem{lemma}[theorem]{Lemma}
\newtheorem{proposition}[theorem]{Proposition}
\newtheorem{corollary}[theorem]{Corollary}

\theoremstyle{definition}
\newtheorem{definition}[theorem]{Definition}
\newtheorem{example}[theorem]{Example}

\theoremstyle{remark}
\newtheorem{remark}[theorem]{Remark}

% ----------------------------------------------------------------------
% Commands
% ----------------------------------------------------------------------
\newcommand{\C}{\mathbb{C}}
\newcommand{\R}{\mathbb{R}}
\newcommand{\Z}{\mathbb{Z}}
\newcommand{\Q}{\mathbb{Q}}
\newcommand{\PP}{\mathbb{P}}
\newcommand{\AAA}{\mathbb{A}}
\newcommand{\OO}{\mathcal{O}}
\newcommand{\Xan}{X^{\mathrm{an}}}
\newcommand{\Fan}{\mathcal{F}^{\mathrm{an}}}
\newcommand{\dd}{\mathrm{d}}
\newcommand{\ii}{\mathrm{i}}
\newcommand{\dbar}{\bar\partial}
\DeclareMathOperator{\Spec}{Spec}
\DeclareMathOperator{\Pic}{Pic}
\DeclareMathOperator{\Coh}{Coh}
\DeclareMathOperator{\Hom}{Hom}
\DeclareMathOperator{\divergence}{div}
\DeclareMathOperator{\id}{id}

% ----------------------------------------------------------------------
% Document information
% ----------------------------------------------------------------------
\title[Analytification, GAGA, and Lefschetz Theory]
{Notes on Analytification, GAGA, K\"ahler Identities,\\
  and Lefschetz Theory}
\author{Jhune Yim}
\email{jhune\_yim@protonmail.com}
\date{12 November 2026}
\subjclass[2020]{Primary 32Q15; Secondary 14C30, 14F25, 14J32, 32C35}
\keywords{Analytification, GAGA, K\"ahler identities, Hodge decomposition,
  Hard Lefschetz theorem, quintic threefold}

\begin{document}
\raggedbottom

\begin{abstract}
Analytification gives a smooth projective complex variety a compact K\"ahler
manifold structure. Starting with the analytification functor and GAGA, we
derive the K\"ahler commutation identities in holomorphic normal coordinates
and use $\Delta_{\dd}=2\Delta_{\dbar}$ to obtain the Hodge decomposition
from the elliptic Hodge theorem. The Lefschetz operators on harmonic forms
give a finite-dimensional $\mathfrak{sl}_2$ representation with
$L^*=\Lambda$. Its weight strings yield Hard Lefschetz and the primitive
decomposition. For a smooth quintic threefold we compute the Hodge diamond
and exhibit the index of the integral Lefschetz map in degree two.
\end{abstract}

\maketitle

\setcounter{tocdepth}{1}
\tableofcontents

\section{Introduction}
\label{sec:intro}

Over $\C$, polynomial equations define schemes and give their sets of complex
points the structure of analytic spaces. The analytification functor relates
these constructions. When the scheme is proper, GAGA compares coherent
algebraic and analytic sheaves and their cohomology. Smooth projective schemes
have a second feature: restricting the Fubini--Study form gives their
analytifications a K\"ahler metric.

We use these facts in two ways. GAGA identifies the algebraic groups
$H^q(X,\Omega^p_{X/\C})$ with the corresponding analytic coherent cohomology
groups. On the compact K\"ahler manifold, the K\"ahler identities identify
the harmonic spaces that realize those groups as Hodge summands. The same
identities make the Lefschetz operators preserve harmonic forms. This puts the
proof of Hard Lefschetz within the linear algebra of an $\mathfrak{sl}_2$ action
\cite{griffithsharris1978,voisin2002hodge,huybrechts2005complex}.

The smooth quintic illustrates the two comparisons. Algebraic adjunction
determines $h^{3,0}$, the hyperplane class determines the degree-four Hodge
summand, and the tangent bundle gives the Euler characteristic. The companion
notes \cite{hodgenotes} use the resulting degree-four calculation to
represent the generator of $H^4$ by a curve: rationally by a plane section,
and integrally by a line.

\subsection{Conventions}
\label{subsec:conventions}

We allow non-Hausdorff complex analytic spaces, as required for the
analytification of nonseparated schemes. We use \emph{complex manifold}
for a Hausdorff space with holomorphic manifold charts and a countable base.
The notation for cohomology depends on the coefficient object.
\begin{enumerate}
  \item $H^i(X,\mathcal F)$ denotes Zariski sheaf cohomology on the scheme $X$.
  \item $H^i(\Xan,\mathcal G)$ denotes analytic sheaf cohomology.
  \item $H^k(\Xan,A)$, for $A\in\{\Z,\Q,\R,\C\}$, denotes singular
  cohomology of the underlying topological space.
\end{enumerate}
From \cref{sec:kahler-identities} onward, the ambient object is a K\"ahler
manifold, again denoted by $X$. When it is the analytification of an algebraic
scheme, this is stated explicitly. We put
\[
b_k=\dim_{\R}H^k(X,\R),
\qquad
h^{p,q}=\dim_{\C}H^{p,q}(X).
\]
\Cref{sec:limits} records the foundational results used without proof.

\section{Complex Algebraic Varieties and Analytification}
\label{sec:analytification}

\subsection{Algebraic varieties over \texorpdfstring{$\C$}{C}}

\begin{definition}[Complex algebraic variety]
\label{def:variety}
A \emph{complex algebraic variety} is an integral, separated scheme of finite
type over $\C$.
\end{definition}

Here \emph{integral} means nonempty, irreducible, and reduced. Finite type
implies that $X$ has an affine open cover by schemes $\Spec A$ with
\[
A\cong\C[x_1,\dots,x_N]/I
\]
for a finitely generated ideal $I$. Separatedness means that the diagonal
\[
\Delta_X:X\longrightarrow X\times_{\C}X
\]
is a closed immersion. It is the algebraic counterpart of Hausdorffness in the
precise sense of \cref{prop:sep-proper}, although the Zariski topology itself is
not Hausdorff in positive dimension.

\begin{remark}
\label{rem:conventions-variety}
Terminology varies. Some authors use \emph{variety} for every reduced
finite-type scheme, without irreducibility. Analytification itself requires
neither integrality nor reducedness, so several results below are stated for
arbitrary finite-type $\C$-schemes.
\end{remark}

\subsection{The analytification functor}

Let
\[
X=\Spec\bigl(\C[z_1,\dots,z_N]/I\bigr).
\]
Its complex points are
\[
X(\C)=\{a\in\C^N:f(a)=0\text{ for every }f\in I\},
\]
with the subspace topology from $\C^N$, and with structure sheaf
\[
\OO_{\Xan}
=\bigl(\OO_{\C^N}/I\OO_{\C^N}\bigr)\big|_{X(\C)}.
\]
If $I$ is not radical, this sheaf retains nilpotents. Analytification does not
silently replace a scheme by its reduction.

\begin{theorem}[Analytification, {\cite[Exp.~XII, Thm.\ and Def.~1.1]{sga1}}]
\label{thm:analytification}
There is a functor
\[
(\cdot)^{\mathrm{an}}:
\{\text{schemes locally of finite type over }\C\}
\longrightarrow
\{\text{complex analytic spaces}\}
\]
assigning to $X$ a complex analytic space $\Xan$ and a morphism of locally
ringed spaces over $\C$
\[
\rho_X:\Xan\longrightarrow X,
\]
characterized by the property that for every complex analytic space $Z$,
composition with $\rho_X$ induces a natural bijection
\[
\Hom_{\mathrm{An}}(Z,\Xan)
\xrightarrow{\sim}
\Hom_{\mathrm{LRS}/\C}(Z,X).
\]
If $X$ is of finite type over $\C$, the points of $\Xan$ are naturally
identified with the closed points of $X$, equivalently with $X(\C)$. Locally,
when $X$ is closed in $\AAA^N_{\C}$, its analytic topology and structure sheaf
are the ones described above.
\end{theorem}

The identification of points uses the Nullstellensatz: every closed point of a
finite-type $\C$-scheme has residue field $\C$.

\begin{remark}
\label{rem:topologies}
On the common set $X(\C)$, the analytic topology is finer than the topology
induced from the Zariski topology. If $X$ is irreducible and
positive-dimensional, it is strictly finer. A proper Zariski-closed subset
has smaller dimension than $X$ at each of its points, so it has empty
analytic interior, and every nonempty Zariski-open subset is analytically
dense. A small coordinate ball around a smooth point of the reduction is
analytically open and not dense, hence not Zariski-open.
\end{remark}

\begin{proposition}
\label{prop:sep-proper}
Let $X$ be of finite type over $\C$.
\begin{enumerate}
  \item $X$ is separated over $\C$ if and only if $\Xan$ is Hausdorff.
  \item $X$ is proper over $\C$ if and only if $\Xan$ is compact and
  Hausdorff.
\end{enumerate}
\end{proposition}

See \cite[Exp.~XII, Prop.~3.1(viii) and Prop.~3.2(v)]{sga1}.

\begin{proposition}
\label{prop:smooth-manifold}
Let $X$ be a finite-type $\C$-scheme, smooth of pure dimension $d$. Then
$\Xan$ has local holomorphic manifold charts of dimension $d$. If $X$ is
separated, $\Xan$ is a complex manifold. If $X$ is projective, $\Xan$ is
also compact.
\end{proposition}

\begin{proof}
In an affine embedding into $\AAA^N_{\C}$, smoothness gives $N-d$
local defining equations with independent differentials at a given closed
point. The holomorphic implicit-function theorem gives the charts. For
separated $X$, \cref{prop:sep-proper} gives Hausdorffness. A finite affine
open cover gives a countable base, since each analytic affine chart has
the subspace topology from a complex affine space. For projective $X$,
the analytification is closed in $\PP^N(\C)$ for some $N$, so it is compact.
\end{proof}

Separatedness cannot be dropped from the second statement. Gluing two copies
of $\AAA^1_{\C}$ along the punctured affine line by the identity gives a
smooth scheme of pure dimension one. In the analytification
every neighborhood of either origin meets the common punctured line in
a punctured disk around zero. The two origins therefore have no disjoint
neighborhoods.

\begin{example}[An affine plane curve]
\label{ex:curve}
Let
\[
X=\Spec\bigl(\C[x,y]/(y^2-x^3+x)\bigr).
\]
Then $\Xan\subset\C^2$ is the analytic curve $y^2=x^3-x$. The restriction of
$e^x$ is holomorphic. If it were regular, it would belong to the invariant
part of $\C[x,y]/(y^2-x^3+x)$ under $y\mapsto-y$, namely $\C[x]$. Since the
projection to the $x$-line is surjective, this would make $e^x$ a polynomial
on $\C$. Thus the coherent cohomology comparison in degree zero can fail
for affine schemes.
\end{example}

\section{Serre's GAGA Principle}
\label{sec:gaga}

For a coherent algebraic sheaf $\mathcal F$, define its analytification by
\[
\Fan=\rho_X^*\mathcal F
=\OO_{\Xan}\otimes_{\rho_X^{-1}\OO_X}\rho_X^{-1}\mathcal F.
\]

\begin{theorem}[GAGA]
\label{thm:gaga}
Let $X$ be a proper scheme over $\C$. Analytification induces an equivalence
\[
\Coh(X)\xrightarrow{\sim}\Coh(\Xan),
\]
and for every coherent sheaf $\mathcal F$ on $X$ and every $i\geq0$, the
comparison morphism
\[
H^i(X,\mathcal F)\longrightarrow H^i(\Xan,\Fan)
\]
is an isomorphism of finite-dimensional $\C$-vector spaces.
\end{theorem}

Serre proved the comparison for projective varieties
\cite[\S12, Thms.~1--3]{serre1956gaga}. For proper schemes, Expos\'e XII of
SGA~1 gives the cohomology comparison in Corollary~4.3 and the equivalence
in Theorem~4.4 \cite{sga1}.

\begin{lemma}[Faithful flatness of analytification]
\label{lem:flat}
For $x\in\Xan$ and $y=\rho_X(x)$, the local homomorphism
\[
A=\OO_{X,y}\longrightarrow B=\OO_{\Xan,x}
\]
is faithfully flat. Consequently, analytification is exact on coherent
sheaves, and local freeness of a coherent sheaf, including its rank, may be
tested after analytification.
\end{lemma}

\begin{proof}
The rings $A$ and $B$ are Noetherian local rings. In the affine model at
$a\in\C^N$, the polynomial and convergent power series presentations both
complete to $\C[[z-a]]/(I)$. Thus the natural map on completions
$\widehat A\to\widehat B$ is an isomorphism
\cite[Exp.~XII, Thm.\ and Def.~1.1]{sga1}. The composite
$A\to B\to\widehat B\cong\widehat A$ is flat, and
$B\to\widehat B$ is faithfully flat. Hence $A\to B$ is flat. As a flat
local map of local rings it is faithfully flat.

Exactness follows by tensoring stalkwise with $B$. A coherent module is
finitely presented, so its finite projectivity descends under faithfully
flat base change. Finite projective modules over a local ring are free,
with rank preserved by base change. The locally free locus of a coherent
sheaf is open, and every nonempty closed subset of a finite-type
$\C$-scheme contains a closed point. Thus it suffices to test local freeness
at the algebraic stalks corresponding to analytic points.
\end{proof}

\begin{corollary}
\label{cor:gaga}
Let $X$ be proper over $\C$.
\begin{enumerate}
  \item Every coherent analytic sheaf $\mathcal G$ on $\Xan$ admits an
  algebraization $(\mathcal F,\phi)$, where $\mathcal F$ is coherent algebraic
  and $\phi:\Fan\xrightarrow{\sim}\mathcal G$. Two such pairs have a unique
  isomorphism compatible with the specified identifications.
  \item Every closed analytic subspace of $\Xan$ is the analytification of a
  unique closed algebraic subscheme of $X$.
  \item Every holomorphic vector bundle on $\Xan$ is the analytification of an
  algebraic vector bundle on $X$.
\end{enumerate}
\end{corollary}

\begin{proof}
For (i), essential surjectivity gives the pair $(\mathcal F,\phi)$. For two
such pairs, full faithfulness lifts the isomorphism
$(\phi')^{-1}\phi$ uniquely to an algebraic isomorphism. For (ii), a closed
analytic subspace is determined by a coherent ideal
$\mathcal I\subseteq\OO_{\Xan}$. By GAGA, both $\mathcal I$ and its inclusion
into $\OO_{\Xan}$ descend uniquely to a coherent sheaf $\mathcal J$ and a
morphism $\mathcal J\to\OO_X$. Exactness and faithfulness of analytification
show that this morphism is injective, so $\mathcal J$ is an ideal sheaf. For
(iii), GAGA algebraizes the coherent locally free sheaf of holomorphic
sections, and \cref{lem:flat} descends its local freeness and rank.
\end{proof}

Applying (iii) in rank one, and using full faithfulness for isomorphisms,
gives the comparison $\Pic(X)\cong\Pic(\Xan)$.

\begin{remark}[Chow's theorem]
\label{rem:chow}
A closed analytic subset of $\PP^N(\C)$ is algebraic \cite{chow1949}. This is
the reduced form of \cref{cor:gaga}\,(ii). In the affine line the zero set
of $\sin z$ is an infinite discrete closed analytic subset, whereas every
proper Zariski-closed subset is finite.
\end{remark}

\section{Smooth Projective Varieties as K\"ahler Manifolds}
\label{sec:kahler}

Let $(M,J)$ be a complex manifold. A Riemannian metric $g$ is \emph{Hermitian}
if $g(Ju,Jv)=g(u,v)$ for all real tangent vectors $u,v$.

\begin{definition}[K\"ahler manifold]
\label{def:kahler}
The fundamental form of a Hermitian manifold $(M,J,g)$ is
\[
\omega(u,v)=g(Ju,v).
\]
It is a real differential form of type $(1,1)$. The metric is \emph{K\"ahler}
if $\dd\omega=0$. In that case $\omega$ is called the K\"ahler form.
\end{definition}

On the affine chart $Z_0\neq0$ of $\PP^N(\C)$, with $z_j=Z_j/Z_0$, put
\[
\omega_{\mathrm{FS}}
=\frac{\ii}{2\pi}\,\partial\dbar
\log\left(1+\sum_{j=1}^N|z_j|^2\right).
\]
The local potentials differ on overlaps by the logarithm of the squared norm
of a nowhere-vanishing holomorphic function, so these forms glue to a global
positive closed real $(1,1)$-form. With this normalization,
$\int_{\PP^1}\omega_{\mathrm{FS}}=1$.

\begin{theorem}
\label{thm:proj-kahler}
Let $X$ be a smooth projective $\C$-scheme of pure dimension. Then $\Xan$ is a
compact K\"ahler manifold. Its K\"ahler class may be chosen in the image of
\[
H^2(\Xan,\Z)\longrightarrow H^2(\Xan,\R).
\]
\end{theorem}

\begin{proof}
Choose a closed immersion $X\hookrightarrow\PP^N_{\C}$. Then $\Xan$ is a compact
complex submanifold of $\PP^N(\C)$. The restriction of
$\omega_{\mathrm{FS}}$ is closed, real, of type $(1,1)$, and positive on every
complex tangent space, hence K\"ahler. Its cohomology class is the restriction of
the integral hyperplane class. Indeed, for $N\geq1$, the integral of a
degree-two class over a projective line determines its coefficient in
the hyperplane basis of $H^2(\PP^N(\C),\Z)$.
\end{proof}

\begin{remark}[Kodaira embedding]
\label{rem:kodaira}
A compact K\"ahler manifold is projective if and only if it has a K\"ahler
class in the image of $H^2(M,\Z)\to H^2(M,\R)$, equivalently if it possesses
a positive holomorphic line bundle \cite[\S5.3]{huybrechts2005complex}.
A very general complex torus of dimension at least two has no such class.
\end{remark}

\section{The K\"ahler Identities and the Laplacians}
\label{sec:kahler-identities}

Let $(X,\omega)$ be a K\"ahler manifold of complex dimension $n$.
The identities in this section are local. Write
\[
A^k(X,\C)=\bigoplus_{p+q=k}A^{p,q}(X),
\qquad
\dd=\partial+\dbar.
\]
Let $\partial^*$, $\dbar^*$, and $\dd^*$ be the formal adjoints with respect to
the $L^2$-pairing on compactly supported forms induced by the Hermitian metric
and its Riemannian volume form. Taking formal adjoints is additive, so
\[
\dd^*=\partial^*+\dbar^*.
\]
Set
\[
L\alpha=\omega\wedge\alpha,
\qquad
\Lambda=L^*.
\]
We use ordinary commutators $[A,B]=AB-BA$ when one operator is even and the
anticommutator $[D_1,D_2]_+=D_1D_2+D_2D_1$ for two odd operators.

\subsection{K\"ahler normal coordinates}

Extend the real metric $g$ complex bilinearly to $g_{\C}$. In holomorphic
coordinates write
$g_{j\bar k}=g_{\C}(\partial/\partial z_j,\partial/\partial\bar z_k)$.
With this convention $\omega=\ii\sum_{j,k}g_{j\bar k}\,\dd z_j\wedge
\dd\bar z_k$, and the coordinate coframe is unitary at a point where
$(g_{j\bar k})$ is the identity.

\begin{lemma}[K\"ahler normal coordinates]
\label{lem:normal-coords}
A Hermitian metric $g$ on a complex manifold $X$ is K\"ahler if and only if, for
every $x\in X$, there exist holomorphic coordinates $(z_1,\dots,z_n)$ centered
at $x$ such that
\[
g_{j\bar k}(z)=\delta_{jk}+O(|z|^2).
\]
\end{lemma}

\begin{proof}
If such coordinates exist at every point, the first derivatives of
$g_{j\bar k}$ vanish at each centre. Hence $\dd\omega$ vanishes at every
point.

Conversely, fix $x\in X$ and assume $\dd\omega=0$. After a complex-linear
coordinate change, arrange $g_{j\bar k}(0)=\delta_{jk}$. Write
\[
g_{j\bar k}(z)
=\delta_{jk}
+\sum_l a_{jkl}z_l
+\sum_l\overline{a_{kjl}}\,\bar z_l
+O(|z|^2).
\]
The equation $\partial\omega=0$ gives
\[
\frac{\partial g_{j\bar k}}{\partial z_l}(0)
=
\frac{\partial g_{l\bar k}}{\partial z_j}(0),
\qquad
\text{hence}
\qquad
a_{jkl}=a_{lkj}.
\]
Make a holomorphic quadratic change of coordinates
\[
z_j=w_j+\frac12\sum_{k,l}c_{jkl}w_kw_l,
\qquad c_{jkl}=c_{jlk}.
\]
In the transformed coefficient $g'_{l\bar k}$, the coefficient of $w_m$
is $a_{lkm}+c_{kml}$. The second term comes from differentiating $z$ with
respect to $w_l$ in the metric transformation formula. Choosing
\[
c_{kml}=-a_{lkm}
\]
kills these terms. The required symmetry of $c$ follows from the
K\"ahler symmetry $a_{lkm}=a_{mkl}$. Hermitian symmetry simultaneously removes
the antiholomorphic linear terms. Thus
$g_{j\bar k}(w)=\delta_{jk}+O(|w|^2)$ at $x$.
\end{proof}

In the underlying real coordinates of a holomorphic normal chart, the first
derivatives of the real metric vanish at its centre. The Levi-Civita
connection coefficients vanish there and $J$ has constant coordinate
components. Applying this at each point gives $\nabla J=0$. Since $\nabla g=0$,
it follows that $\nabla\omega=0$ and the induced connection on forms
preserves bidegree.

Fix $x\in X$ and take the holomorphic normal coframe $\theta^j=\dd z_j$,
unitary at $x$. Then
\[
\omega_x=\ii\sum_{j=1}^n\theta^j\wedge\bar\theta^j.
\]
Let
\[
\varepsilon_j\alpha=\theta^j\wedge\alpha,
\qquad
\bar\varepsilon_j\alpha=\bar\theta^j\wedge\alpha,
\]
and let $\iota_j,\bar\iota_j$ denote contraction by the dual $(1,0)$- and
$(0,1)$-vectors. The canonical anticommutation relations are
\[
[\iota_j,\varepsilon_k]_+=\delta_{jk},
\qquad
[\bar\iota_j,\bar\varepsilon_k]_+=\delta_{jk},
\]
with all mixed anticommutators equal to zero. At $x$,
\[
L=\ii\sum_j\varepsilon_j\bar\varepsilon_j,
\qquad
\Lambda=-\ii\sum_j\bar\iota_j\iota_j.
\]

\begin{lemma}
\label{lem:lambda-eps}
At $x$,
\[
[\Lambda,\varepsilon_k]=-\ii\bar\iota_k,
\qquad
[\Lambda,\bar\varepsilon_k]=\ii\iota_k.
\]
\end{lemma}

\begin{proof}
Terms with index $j\neq k$ commute with $\varepsilon_k$ and
$\bar\varepsilon_k$. For $j=k$,
\[
\bar\iota_k\iota_k\varepsilon_k
=\bar\iota_k(1-\varepsilon_k\iota_k)
=\bar\iota_k+\varepsilon_k\bar\iota_k\iota_k,
\]
so $[\bar\iota_k\iota_k,\varepsilon_k]=\bar\iota_k$. Similarly,
\[
\bar\iota_k\iota_k\bar\varepsilon_k
=-\bar\iota_k\bar\varepsilon_k\iota_k
=-\iota_k+\bar\varepsilon_k\bar\iota_k\iota_k,
\]
so $[\bar\iota_k\iota_k,\bar\varepsilon_k]=-\iota_k$. Multiplication by
$-\ii$ proves the result.
\end{proof}

\begin{lemma}
\label{lem:L-Lambda}
The operator $[L,\Lambda]$ acts on $A^k(X,\C)$ as multiplication by $k-n$.
\end{lemma}

\begin{proof}
Put $N_j=\varepsilon_j\iota_j$ and
$\bar N_j=\bar\varepsilon_j\bar\iota_j$. On a monomial
$\theta^I\wedge\bar\theta^J$, these are the number operators for the occurrence
of $j$ in $I$ and $J$. Since
$\iota_j\varepsilon_j=1-N_j$ and
$\bar\iota_j\bar\varepsilon_j=1-\bar N_j$, a direct calculation gives
\[
[L,\Lambda]
=\sum_j\left(N_j+\bar N_j-1\right).
\]
On a $(p,q)$-form, the right-hand side is $p+q-n=k-n$.
\end{proof}

\begin{lemma}[Local formulas for the Dolbeault operators]
\label{lem:frame-operators}
In a holomorphic normal frame, write
$\nabla_j=\nabla_{\partial/\partial z_j}$ and
$\nabla_{\bar j}=\nabla_{\partial/\partial\bar z_j}$.
At its centre,
\[
\partial=\sum_j\varepsilon_j\nabla_j,
\qquad
\dbar=\sum_j\bar\varepsilon_j\nabla_{\bar j},
\]
and
\[
\partial^*=-\sum_j\iota_j\nabla_{\bar j},
\qquad
\dbar^*=-\sum_j\bar\iota_j\nabla_j.
\]
\end{lemma}

\begin{proof}
The K\"ahler Levi-Civita connection is torsion-free and preserves type.
Exterior differentiation is therefore the alternation of $\nabla$, and its
$(1,0)$- and $(0,1)$-parts give the first two formulas, which hold throughout
the chart.

For the adjoints, write $\langle\cdot,\cdot\rangle$ for the pointwise
Hermitian product on forms, linear in the first entry. Since $\nabla$ is real
and metric,
\begin{equation}
\label{eq:metric-leibniz}
v\langle\alpha,\gamma\rangle
=\langle\nabla_v\alpha,\gamma\rangle
+\langle\alpha,\nabla_{\bar v}\gamma\rangle
\end{equation}
for every complex vector field $v$. The form $g_{\C}$ pairs $T^{1,0}$
perfectly with $T^{0,1}$, so each covector $\bar\theta$ of type $(0,1)$ has a
dual vector $\bar\theta^{\sharp}\in T^{1,0}$ with
$g_{\C}(\bar\theta^{\sharp},\cdot)=\bar\theta$. Let $\nabla^{0,1}\beta$ be the
component of $\nabla\beta$ in $(T^{0,1})^{\vee}\otimes\bigwedge^{\bullet}$,
and define a first-order operator of bidegree $(-1,0)$ by
\[
D\beta=-c\bigl(\nabla^{0,1}\beta\bigr),
\qquad
c(\bar\theta\otimes\gamma)=\iota_{\bar\theta^{\sharp}}\gamma.
\]
The definition is global and involves no coordinates. At the centre of a
normal chart $(\dd\bar z_j)^{\sharp}=\partial/\partial z_j$, so
$D\beta=-\sum_j\iota_j\nabla_{\bar j}\beta$ there.

Fix forms $\alpha$ and $\beta$, and let $\xi$ be the vector field of type
$(1,0)$ determined by
$g_{\C}(\xi,\bar w)=\langle\alpha,\iota_w\beta\rangle$ for $w\in T^{1,0}$.
We claim that
\begin{equation}
\label{eq:adjoint-divergence}
\langle\partial\alpha,\beta\rangle-\langle\alpha,D\beta\rangle
=\divergence\xi
\end{equation}
at every point. Both sides are globally defined functions, so it suffices to
check the identity at the centre $x$ of a normal chart. The coframe is
unitary at $x$, so $\varepsilon_j$ has pointwise adjoint $\iota_j$ there.
The connection coefficients vanish at $x$, so $\nabla_{\bar j}$ commutes
with $\iota_j$ at $x$. Using \eqref{eq:metric-leibniz} with
$v=\partial/\partial z_j$, we get at $x$
\[
\langle\partial\alpha,\beta\rangle
=\sum_j\langle\nabla_j\alpha,\iota_j\beta\rangle
=\sum_j\frac{\partial}{\partial z_j}\langle\alpha,\iota_j\beta\rangle
-\sum_j\langle\alpha,\iota_j\nabla_{\bar j}\beta\rangle.
\]
Since $D\beta=-\sum_j\iota_j\nabla_{\bar j}\beta$ at $x$, the last term is
$+\langle\alpha,D\beta\rangle$. The components of $\xi$ are
$\xi^k=\sum_l g^{\bar lk}\langle\alpha,\iota_l\beta\rangle$, where
$(g^{\bar lk})$ is the inverse matrix of $(g_{k\bar l})$. The first
derivatives of the metric vanish at $x$, so
$\divergence\xi=\sum_j\partial\xi^j/\partial z_j
=\sum_j\partial\langle\alpha,\iota_j\beta\rangle/\partial z_j$ at $x$.
This proves \eqref{eq:adjoint-divergence}.

If $\alpha$ or $\beta$ has compact support, so does $\xi$, and the divergence
theorem applied to the real and imaginary parts of $\xi$ shows that
$\divergence\xi$ has integral zero. Integrating
\eqref{eq:adjoint-divergence} gives $\partial^*=D$, which is the stated
formula at the centre. The formula for $\dbar^*$ follows by complex
conjugation.
\end{proof}

\begin{theorem}[K\"ahler identities]
\label{thm:kahler-identities}
On a K\"ahler manifold,
\begin{align}
[\Lambda,\partial]&=\ii\dbar^*,
&
[\Lambda,\dbar]&=-\ii\partial^*,
\label{eq:kahler-first}\\
[L,\partial^*]&=\ii\dbar,
&
[L,\dbar^*]&=-\ii\partial,
\label{eq:kahler-second}
\end{align}
and
\[
[L,\partial]=[L,\dbar]=0,
\qquad
[\Lambda,\partial^*]=[\Lambda,\dbar^*]=0.
\]
\end{theorem}

\begin{proof}
At a fixed point, $\nabla\omega=0$ implies that $\Lambda$ commutes with the
covariant derivatives. By \cref{lem:lambda-eps,lem:frame-operators},
\[
[\Lambda,\partial]
=\sum_j[\Lambda,\varepsilon_j]\nabla_j
=-\ii\sum_j\bar\iota_j\nabla_j
=\ii\dbar^*,
\]
and similarly
\[
[\Lambda,\dbar]
=\ii\sum_j\iota_j\nabla_{\bar j}
=-\ii\partial^*.
\]
Taking formal adjoints gives \eqref{eq:kahler-second}. Finally,
$\partial\omega=\dbar\omega=0$, so wedging by $\omega$ commutes with
$\partial$ and $\dbar$. Taking adjoints yields the remaining identities.
\end{proof}

On a general Hermitian manifold the pointwise linear-algebra identity
\cref{lem:L-Lambda} remains true, but the differential K\"ahler identities
acquire torsion terms and the equality of Laplacians below generally fails.

\subsection{Equality of the Laplacians}

Define
\[
\Delta_\partial=\partial\partial^*+\partial^*\partial,
\quad
\Delta_{\dbar}=\dbar\dbar^*+\dbar^*\dbar,
\quad
\Delta_\dd=\dd\dd^*+\dd^*\dd.
\]

\begin{lemma}
\label{lem:mixed-terms}
The mixed anticommutators vanish:
\[
[\partial,\dbar^*]_+=0,
\qquad
[\dbar,\partial^*]_+=0.
\]
\end{lemma}

\begin{proof}
Since $\dbar^*=-\ii[\Lambda,\partial]$,
\[
\partial\dbar^*+\dbar^*\partial
=-\ii\bigl(
\partial\Lambda\partial-\partial^2\Lambda
+\Lambda\partial^2-\partial\Lambda\partial
\bigr)=0.
\]
The second identity follows similarly from
$\partial^*=\ii[\Lambda,\dbar]$ and $\dbar^2=0$.
\end{proof}

\begin{theorem}
\label{thm:laplacian-equality}
On a K\"ahler manifold,
\[
\Delta_\partial=\Delta_{\dbar},
\qquad
\Delta_\dd=2\Delta_{\dbar}=2\Delta_\partial.
\]
\end{theorem}

\begin{proof}
Using $\partial^*=\ii[\Lambda,\dbar]$ and
$\partial\dbar+\dbar\partial=0$,
\[
\begin{aligned}
\Delta_\partial
&=\ii\bigl(\partial[\Lambda,\dbar]+[\Lambda,\dbar]\partial\bigr)\\
&=\ii\bigl((\partial\Lambda-\Lambda\partial)\dbar
+\dbar(\partial\Lambda-\Lambda\partial)\bigr).
\end{aligned}
\]
By $[\Lambda,\partial]=\ii\dbar^*$,
$\partial\Lambda-\Lambda\partial=-\ii\dbar^*$, and therefore
\[
\Delta_\partial
=\dbar^*\dbar+\dbar\dbar^*
=\Delta_{\dbar}.
\]
Expanding $\dd=\partial+\dbar$ gives
\[
\Delta_\dd
=\Delta_\partial+\Delta_{\dbar}
+[\partial,\dbar^*]_++[\dbar,\partial^*]_+.
\]
The mixed terms vanish by \cref{lem:mixed-terms}, proving the result.
\end{proof}

\begin{corollary}
\label{cor:harmonic-agree}
On smooth complex forms, $\ker\Delta_\dd=\ker\Delta_\partial
=\ker\Delta_{\dbar}$. Moreover, $\Delta_\dd$ preserves bidegree.
\end{corollary}

\section{Hodge Decomposition}
\label{sec:hodge}

From now on, $X$ is compact K\"ahler of complex dimension $n$.
We use the elliptic Hodge theorem for the de Rham and Dolbeault complexes.
Each class in either complex has a unique harmonic representative for the
corresponding Laplacian, and the harmonic spaces are finite-dimensional.

\begin{theorem}[Hodge decomposition]
\label{thm:hodge}
For every $k$,
\[
H^k(X,\C)=\bigoplus_{p+q=k}H^{p,q}(X),
\qquad
H^{p,q}(X)\cong H^q(X,\Omega_X^p).
\]
The subspace $H^{p,q}(X)$ consists intrinsically of those cohomology classes
represented by a closed form of type $(p,q)$. Consequently the decomposition
is independent of the chosen K\"ahler metric and is functorial for
holomorphic maps.
Complex conjugation gives an antilinear isomorphism
\[
H^{p,q}(X)\xrightarrow{\sim}H^{q,p}(X),
\]
so $h^{p,q}=h^{q,p}$.
\end{theorem}

\begin{proof}
By \cref{thm:laplacian-equality}, the operator $\Delta_\dd$ preserves
bidegree and its kernel on $A^{p,q}(X)$ is the kernel of $\Delta_{\dbar}$
on that space. The Hodge theorem identifies the degree-$k$ de Rham
cohomology with the direct sum of these kernels for $p+q=k$.
The Hodge theorem for $\dbar$ and the Dolbeault isomorphism identify
the $(p,q)$ kernel with $H^q(X,\Omega_X^p)$.

For a closed form of pure type, its $L^2$ projection onto harmonic forms
represents the same de Rham class. The projection preserves type because
the harmonic space splits orthogonally by bidegree. Conversely a harmonic
form of pure type is a closed representative of that type. This gives
the intrinsic description. Pullback by a holomorphic map preserves
closed forms and their type, and complex conjugation commutes with the
real Laplacian $\Delta_\dd$ while exchanging $(p,q)$ and $(q,p)$.
\end{proof}

Wedge products of closed forms give the corresponding bidegree rule for
cup products,
$H^{p,q}\smile H^{r,s}\subseteq H^{p+r,q+s}$.
See \cite[\S6.1]{voisin2002hodge} and
\cite[\S3.2]{huybrechts2005complex} for the K\"ahler Hodge theorem.

The decomposition defines a pure integral Hodge structure of weight $k$ on the
free lattice
\[
H^k(X,\Z)_{\mathrm{free}}
=H^k(X,\Z)/H^k(X,\Z)_{\mathrm{tors}},
\]
whose rationalization is the pure rational Hodge structure on $H^k(X,\Q)$.
Its Hodge filtration is
\[
F^pH^k(X,\C)=\bigoplus_{r\geq p}H^{r,k-r}(X),
\qquad
H^k(X,\C)=F^p\oplus\overline{F^{k-p+1}}.
\]
This construction applies to every compact K\"ahler manifold. The role of
the K\"ahler class in polarizing the Hodge structure is discussed in
\cref{rem:polarization}.

\subsection{Algebraic computation of Hodge numbers}
\label{subsec:algebraicity}

\begin{corollary}
\label{cor:algebraic-hodge}
Let $X$ be a smooth projective $\C$-scheme of pure dimension. Then
\[
H^q\bigl(X,\Omega^p_{X/\C}\bigr)
\cong
H^q\bigl(\Xan,\Omega^p_{\Xan}\bigr)
\cong
H^{p,q}(\Xan).
\]
Consequently,
\[
h^{p,q}(\Xan)
=\dim_{\C}H^q\bigl(X,\Omega^p_{X/\C}\bigr).
\]
\end{corollary}

\begin{proof}
In a local affine embedding, the conormal presentation of
$\Omega^1_{X/\C}$ analytifies to the presentation of holomorphic
one-forms by the same differentials of local equations. Exterior powers
give $(\Omega^p_{X/\C})^{\mathrm{an}}\cong\Omega^p_{\Xan}$.
The sheaf $\Omega^p_{X/\C}$ is coherent, so GAGA gives the first
isomorphism. \Cref{thm:hodge} gives the second.
\end{proof}

In the route adopted here, this comparison permits algebraic calculations with
$\Omega^p_{X/\C}$ to be interpreted as statements about the Hodge decomposition
of the analytic manifold $\Xan$.

\begin{corollary}
\label{cor:odd-betti}
Every odd Betti number of a compact K\"ahler manifold is even.
\end{corollary}

\begin{proof}
If $p+q$ is odd, then $p\neq q$. The summands in
$\sum_{p+q=2r+1}h^{p,q}$ therefore occur in distinct conjugate pairs of equal
dimension.
\end{proof}

\begin{example}[Projective space]
\label{ex:pn}
For $\PP^n(\C)$,
\[
h^{p,q}=0\quad(p\neq q),
\qquad
h^{p,p}=1\quad(0\leq p\leq n),
\]
and
\[
H^*(\PP^n(\C),\Z)\cong\Z[h]/(h^{n+1}),
\qquad \deg h=2.
\]
The ring structure is classical. The hyperplane class $h$ has type $(1,1)$,
so $h^p$ spans $H^{p,p}$, which gives the Hodge numbers.
\end{example}

\begin{example}[The Hopf surface]
\label{ex:hopf}
Let
\[
M=(\C^2\setminus\{0\})/\langle z\mapsto2z\rangle.
\]
The map with coordinates $\log|z|$ modulo $\log 2$ and $z/|z|$
gives a diffeomorphism $M\cong S^1\times S^3$ \cite[\S V.18]{bhpv2004}.
In particular, $b_1(M)=1$ excludes a K\"ahler metric by
\cref{cor:odd-betti}.

Here the notation $h^{p,q}$ means $\dim_{\C}H^q(M,\Omega_M^p)$.
A holomorphic $p$-form on $M$ for $p=1,2$ pulls back to a form whose
coefficients extend across zero by Hartogs' theorem. Each coefficient
satisfies $2^p f(2z)=f(z)$ by invariance under dilation, so its Taylor
series vanishes. Thus $h^{1,0}=h^{2,0}=0$. The surface Riemann--Roch
formula gives
\[
\chi(M,\OO_M)=\frac{1}{12}\int_M\bigl(c_1(T_M)^2+c_2(T_M)\bigr)=0.
\]
Here $H^2(M,\Z)=0$ and the top Chern number is the Euler characteristic
of $S^1\times S^3$, which is zero
\cite[\S5.1]{huybrechts2005complex}. Serre duality gives $h^{0,2}=0$.
Since $h^{0,0}=1$, it follows that $h^{0,1}=1$. Thus $h^{1,0}=0$ and
$h^{0,1}=1$: Hodge symmetry fails on $M$.
\end{example}

\section{Hard Lefschetz from the
\texorpdfstring{$\mathfrak{sl}_2$}{sl2}-Action}
\label{sec:hard-lefschetz}

Put
\[
\mathcal H=[L,\Lambda].
\]
By \cref{lem:L-Lambda}, $\mathcal H$ acts on $k$-forms by $k-n$. Hence
\[
[\mathcal H,L]=2L,
\qquad
[\mathcal H,\Lambda]=-2\Lambda,
\qquad
[L,\Lambda]=\mathcal H.
\]
Moreover, $L^*=\Lambda$ and $\mathcal H^*=\mathcal H$.

\begin{proposition}
\label{prop:lefschetz-preserves-harmonic}
The operators $L$, $\Lambda$, and $\mathcal H$ commute with $\Delta_\dd$. Thus
they preserve harmonic forms and make $H^*(X,\C)$ a finite-dimensional
Hermitian $\mathfrak{sl}_2(\C)$-module satisfying $L^*=\Lambda$. On cohomology,
$L$ is cup product with $[\omega]$.
\end{proposition}

\begin{proof}
The K\"ahler identities give $[L,\dbar]=0$ and
$[\dbar^*,L]=\ii\partial$. Therefore
\[
[\Delta_{\dbar},L]
=\dbar[\dbar^*,L]+[\dbar^*,L]\dbar
=\ii(\dbar\partial+\partial\dbar)=0.
\]
By \cref{thm:laplacian-equality}, $[\Delta_\dd,L]=0$. Taking adjoints gives
$[\Delta_\dd,\Lambda]=0$, and then $[\Delta_\dd,\mathcal H]=0$. The K\"ahler
form is parallel, hence harmonic, and wedging by it induces cup product with
$[\omega]$ in de Rham cohomology.
\end{proof}

\begin{lemma}[Finite-dimensional Hermitian $\mathfrak{sl}_2$-modules]
\label{lem:sl2-module}
Let $V$ be a finite-dimensional Hermitian vector space with operators $e,f,h$
satisfying
\[
[h,e]=2e,
\qquad
[h,f]=-2f,
\qquad
[e,f]=h,
\]
and $e^*=f$, $h^*=h$. Write $V_a=\ker(h-a\id)$ and $P=\ker f$. Then:
\begin{enumerate}
  \item $V$ is an orthogonal direct sum of irreducible strings
  \[
  \operatorname{span}(w_0,ew_0,\dots,e^mw_0),
  \quad
  fw_0=0,
  \quad
  hw_0=-mw_0,
  \]
  with $m\in\Z_{\geq0}$.
  \item $P\subseteq\bigoplus_{a\leq0}V_a$.
  \item $V=\bigoplus_{j\geq0}e^jP$.
  \item For every $r\in\Z_{\geq0}$, $e^r:V_{-r}\to V_r$ is an isomorphism.
  \item For $a\in\Z_{\leq0}$,
  \[
  P\cap V_a=\ker(e^{-a+1})\cap V_a.
  \]
\end{enumerate}
\end{lemma}

\begin{proof}
Because $h$ is self-adjoint, it is diagonalizable with real eigenvalues. The
operators $e$ and $f$ raise and lower weights by two. Choose a nonzero vector
$w_0$ of minimal weight. Then $fw_0=0$. Write $hw_0=-mw_0$. Induction from
$fe=ef-h$ gives
\[
f e^j w_0=j(m-j+1)e^{j-1}w_0,\qquad j\geq1.
\]
If $\ell$ is maximal with $e^\ell w_0\neq0$, applying this formula to
$e^{\ell+1}w_0=0$ gives $m=\ell\in\Z_{\geq0}$. Thus the span of
$w_0,ew_0,\dots,e^mw_0$ is invariant and has weights
$-m,-m+2,\dots,m$. Any nonzero invariant subspace contains a weight vector
by projecting onto an eigenspace with a polynomial in $h$. Each weight
space in this string is a line. The recurrence and the action of $e$
generate the full string from any of these lines, proving irreducibility.
The orthogonal complement is invariant because $e^*=f$ and $h^*=h$.
Induction on $\dim V$ proves (i).

On each irreducible string, $\ker f$ is the lowest-weight line $\C w_0$,
proving (ii) and (iii). The weight spaces of weights $-r$ and $r$ are either
both zero or are one-dimensional and are carried isomorphically to one another
by $e^r$, which proves (iv). In a string of lowest weight $-m$, a vector
of weight $a=-m+2t\leq0$ lies at level $t$. Applying $e^{1-a}$ would
send it to level $m-t+1$. This exceeds the last level $m$ exactly when
$t=0$, giving (v).
\end{proof}

\begin{theorem}[Hard Lefschetz]
\label{thm:hard-lefschetz}
Let $X$ be compact K\"ahler of complex dimension $n$. For every integer
$0\leq r\leq n$,
\[
L^r:H^{n-r}(X,\C)\xrightarrow{\sim}H^{n+r}(X,\C).
\]
\end{theorem}

\begin{proof}
Apply \cref{lem:sl2-module} with
$e=L$, $f=\Lambda$, and $h=\mathcal H$. Since $\mathcal H$ acts on $H^k$ by
$k-n$, the weight spaces $V_{-r}$ and $V_r$ are respectively
$H^{n-r}(X,\C)$ and $H^{n+r}(X,\C)$. Now use
\cref{lem:sl2-module}\,(iv).
\end{proof}

\begin{definition}[Primitive cohomology]
\label{def:primitive}
On cohomology, $\Lambda$ acts on harmonic representatives. Set
\[
P^k(X)=\ker\bigl(\Lambda:H^k(X,\C)\to H^{k-2}(X,\C)\bigr).
\]
\end{definition}

By \cref{lem:sl2-module}\,(ii), $P^k(X)=0$ for $k>n$.

\begin{proposition}
\label{prop:primitive-alt}
For $k\leq n$,
\[
P^k(X)
=\ker\bigl(L^{n-k+1}:H^k(X,\C)\to H^{2n-k+2}(X,\C)\bigr).
\]
\end{proposition}

\begin{proof}
The group $H^k$ is the weight space of weight $a=k-n\leq0$. Apply
\cref{lem:sl2-module}\,(v), noting that $-a+1=n-k+1$.
\end{proof}

\begin{theorem}[Lefschetz decomposition]
\label{thm:lefschetz-decomp}
For $0\leq k\leq2n$,
\[
H^k(X,\C)=\bigoplus_{j=\max(0,k-n)}^{\lfloor k/2\rfloor}
L^jP^{k-2j}(X).
\]
\end{theorem}

\begin{proof}
Take the degree-$k$ part of the string decomposition in
\cref{lem:sl2-module}. A string beginning in degree $s=k-2j$ has
lowest weight $s-n$ and last term $L^{n-s}$ applied to a primitive
vector. Thus its $j$-th translate occurs for $s\geq0$ and
$0\leq j\leq n-s$, which gives the stated bounds.
\end{proof}

\begin{remark}[Rational coefficients]
\label{rem:rational-lefschetz}
If $[\omega]$ is rational, cup product by $[\omega]$ defines $L$ over
$\Q$. The rational Hard Lefschetz map becomes an isomorphism over $\C$
and is therefore an isomorphism over $\Q$. For $0\leq k\leq n$ put
\[
P^k(X,\Q)=\ker\bigl(L^{n-k+1}:H^k(X,\Q)\to H^{2n-k+2}(X,\Q)\bigr).
\]
Complexification identifies this with the complex primitive space
$P^k(X,\C)=P^k(X)$ by \cref{prop:primitive-alt}. The sum of the rational images
$L^jP^{k-2j}(X,\Q)$ likewise becomes the Lefschetz decomposition over
$\C$, so the direct sum holds over $\Q$. Cup product by $[\omega]$ has
type $(1,1)$, and the primitive kernels inherit their Hodge structures
from cohomology.
\end{remark}

\begin{corollary}
\label{cor:betti-monotone}
For $k<n$, the map
\[
L:H^k(X,\C)\longrightarrow H^{k+2}(X,\C)
\]
is injective. Hence $b_k\leq b_{k+2}$.
\end{corollary}

\begin{proof}
If $L\alpha=0$, then $L^{n-k}\alpha=0$. Hard Lefschetz says that
$L^{n-k}:H^k\to H^{2n-k}$ is an isomorphism, so $\alpha=0$.
\end{proof}

The inequality need not be strict: all nonzero even Betti numbers of
$\PP^n(\C)$ are equal to one.

\section{The Hodge Numbers of a Smooth Quintic Threefold}
\label{sec:quintic}

Let $X=V(F)\subset\PP^4_{\C}$ be a smooth hypersurface of degree five and
put $M=\Xan$. Use the normalized Fubini--Study form on $M$, whose class
is the image of
\[
h=c_1\bigl(\OO_{\PP^4}(1)|_X\bigr)\in H^2(M,\Z).
\]
In this section the Betti and Hodge numbers are those of $M$.
Chern classes of algebraic bundles are taken in topological cohomology
on $M$.

\subsection{Weak Lefschetz and adjunction}

The degree-five Veronese embedding presents $X$ as a hyperplane section
of the embedded $\PP^4$. The weak Lefschetz theorem gives
\[
H^k(\PP^4(\C),\Z)\xrightarrow{\sim}H^k(M,\Z)
\qquad(k<3)
\]
\cite[Ch.~1]{voisin2003hodge2}. In degree zero this gives connectedness.
In degrees one and two it gives $b_1=0$ and
$H^2(M,\Z)=\Z h$. The generator $h$ has type $(1,1)$, so
\[
h^{1,0}=h^{0,1}=h^{2,0}=h^{0,2}=0,\qquad h^{1,1}=1.
\]

Adjunction and $K_{\PP^4}\cong\OO_{\PP^4}(-5)$ give
\[
K_X
\cong\bigl(K_{\PP^4}\otimes\OO_{\PP^4}(5)\bigr)|_X
\cong\OO_X.
\]
Holomorphic functions on the connected compact manifold $M$ are constant,
so GAGA gives $\dim_{\C}H^0(X,\OO_X)=1$. By \cref{cor:algebraic-hodge}
and Hodge symmetry,
\[
h^{3,0}=\dim_{\C}H^0(X,K_X)=1,\qquad h^{0,3}=1.
\]
The same comparison gives $H^1(X,\OO_X)=H^2(X,\OO_X)=0$ from the
$(0,1)$ and $(0,2)$ vanishings. Thus $X$ satisfies the Calabi--Yau
convention used here, namely trivial canonical bundle and
$H^i(X,\OO_X)=0$ for $0<i<3$.

Poincar\'e duality gives $b_5=0$ and $b_6=1$. By rational Hard Lefschetz,
cup product by $h$ identifies $H^2(M,\Q)$ with $H^4(M,\Q)$
(\cref{rem:rational-lefschetz}). Hence
\[
H^4(M,\Q)=\Q h^2.
\]
The generator has type $(2,2)$, so $h^{2,2}=1$ and all other
degree-four Hodge numbers vanish.

\subsection{Chern classes and the Euler characteristic}

The normal bundle of $X$ in $\PP^4$ is $\OO_X(5)$, so
\[
0\longrightarrow T_X
\longrightarrow T_{\PP^4}|_X
\longrightarrow\OO_X(5)
\longrightarrow0.
\]
The Euler sequence gives
\[
c(T_{\PP^4}|_X)=(1+h)^5.
\]
Therefore
\[
c(T_X)=\frac{(1+h)^5}{1+5h}.
\]
Through the term in $h^3$,
\[
(1+5h)^{-1}=1-5h+25h^2-125h^3+\cdots,
\]
and hence
\[
c(T_X)=1+10h^2-40h^3.
\]
Thus
\[
c_1(T_X)=0,
\qquad
c_2(T_X)=10h^2,
\qquad
c_3(T_X)=-40h^3.
\]
Since $\deg X=5$,
\[
\int_M h^3=5.
\]
The top Chern class of the complex tangent bundle is the Euler class
of the underlying real tangent bundle. Its integral is
$\chi_{\mathrm{top}}(M)$
\cite[Cor.~11.12 and \S14]{milnorstasheff1974}. Therefore
\[
\chi_{\mathrm{top}}(M)
=\int_M c_3(T_X)
=-40\int_M h^3
=-200.
\]
On the other hand,
\[
\chi_{\mathrm{top}}(M)
=1-0+1-b_3+1-0+1
=4-b_3,
\]
so
\[
b_3=204.
\]
Finally,
\[
H^3(M,\C)
=H^{3,0}\oplus H^{2,1}\oplus H^{1,2}\oplus H^{0,3}.
\]
Hodge symmetry gives $h^{2,1}=h^{1,2}$, and therefore
\[
204=1+h^{2,1}+h^{1,2}+1
=2+2h^{2,1}.
\]
Consequently,
\[
\boxed{h^{2,1}=h^{1,2}=101}.
\]

\begin{theorem}[Hodge diamond of a smooth quintic threefold]
\label{thm:quintic}
The nonzero Hodge numbers of a smooth quintic threefold are
\[
h^{0,0}=h^{3,3}=1,
\quad
h^{1,1}=h^{2,2}=1,
\quad
h^{3,0}=h^{0,3}=1,
\quad
h^{2,1}=h^{1,2}=101.
\]
Thus its Hodge diamond is
\[
\begin{array}{ccccccc}
&&&1&&&\\[2pt]
&&0&&0&&\\[2pt]
&0&&1&&0&\\[2pt]
1&&101&&101&&1\\[2pt]
&0&&1&&0&\\[2pt]
&&0&&0&&\\[2pt]
&&&1&&&
\end{array}
\]
and
\[
(b_0,\dots,b_6)=(1,0,1,204,1,0,1),
\qquad
\chi_{\mathrm{top}}=-200.
\]
\end{theorem}

\begin{remark}[All smooth quintics are diffeomorphic]
\label{rem:quintic-diffeo}
The computation above applies to each smooth quintic separately, so
\cref{thm:quintic} needs no comparison between different quintics. More is
true: any two smooth quintics are diffeomorphic. The smooth quintic
equations form the complement of the discriminant in $\PP^{125}$. The
projective line through two distinct smooth equations meets the
discriminant in finitely many points, and the complement of these points
in the line is path-connected. Over the smooth locus of equations, the
universal hypersurface is a proper submersion, so Ehresmann's theorem
gives diffeomorphic fibers.
\end{remark}

\begin{remark}[Deformation-theoretic check]
\label{rem:moduli}
Choose a nowhere-vanishing algebraic section of $K_X$. Contraction
identifies $T_X\cong\Omega_X^2$, so $H^1(X,T_X)\cong H^{2,1}(M)$ by
\cref{cor:algebraic-hodge}. The same identification and $h^{2,0}=0$
give $H^0(X,T_X)=0$.

The restriction sequence
\[
0\longrightarrow\OO_{\PP^4}(m-5)\longrightarrow\OO_{\PP^4}(m)
\longrightarrow\OO_X(m)\longrightarrow0
\]
gives $h^0(X,\OO_X(5))=\binom{9}{4}-1=125$ and
$h^0(X,\OO_X(1))=5$. Here we use the vanishing of $H^1$ of line bundles
on $\PP^4$ \cite[Ch.~III, \S5]{hartshorne1977}. With $m=1$, the vanishing of
$H^1(\PP^4,\OO_{\PP^4}(1))$ and $H^2(\PP^4,\OO_{\PP^4}(-4))$ also gives
$H^1(X,\OO_X(1))=0$. The restricted Euler sequence
\[
0\longrightarrow\OO_X\longrightarrow\OO_X(1)^{\oplus5}
\longrightarrow T_{\PP^4}|_X\longrightarrow0
\]
and $H^1(X,\OO_X)=0$ give $h^0(X,T_{\PP^4}|_X)=5\cdot5-1=24$.
Its long exact sequence, with $H^1(X,\OO_X(1))=0$ and $H^2(X,\OO_X)=0$,
gives $H^1(X,T_{\PP^4}|_X)=0$.

The normal sequence therefore gives $h^1(X,T_X)=125-24=101$.
First-order abstract deformations of $X$ are classified by $H^1(X,T_X)$
\cite[Ch.~6]{huybrechts2005complex}. The connecting map of the normal
sequence takes a first-order embedded deformation to its abstract
deformation class \cite[Ch.~3]{sernesi2006}. Its surjectivity shows that
every first-order deformation comes from varying the quintic equation.
\end{remark}

\section{Topological and Hodge-Theoretic Consequences}
\label{sec:consequences}

Let $X$ be nonempty, smooth, and projective over $\C$ of pure dimension $n$,
and put $M=\Xan$. The preceding theorems give the following constraints.
\begin{enumerate}
  \item Every odd Betti number of $M$ is even.
  \item Poincar\'e duality gives $b_k=b_{2n-k}$.
  \item Hard Lefschetz gives $b_k\leq b_{k+2}$ for $k<n$. Together with
  Poincar\'e duality, the even-degree and odd-degree Betti subsequences are
  separately symmetric and unimodal about the middle dimension. For a
  curve $C_g$ of genus $g\geq2$, K\"unneth gives the Betti numbers
  $(1,2g,2,2g,1)$ for $\PP^1\times C_g$, which fail unimodality across
  all degrees.
  \item For every $0\leq k\leq n$, one has $b_{2k}\geq1$. Indeed,
  $\int_M\omega^n>0$, so $[\omega]^n\neq0$, and
  $[\omega]^k[\omega]^{n-k}=[\omega]^n$ forces $[\omega]^k\neq0$.
  Consequently, no smooth projective variety of complex dimension at least two
  is a rational homology sphere. The one-dimensional exception is
  $\PP^1(\C)\cong S^2$.
  \item Hodge symmetry gives $h^{p,q}=h^{q,p}$.
  \item Serre duality and $(\Omega_M^p)^\vee\otimes K_M\cong\Omega_M^{n-p}$
  give $h^{p,q}=h^{n-p,n-q}$.
  \item Each $H^k(M,\Q)$ carries a pure rational Hodge structure. The integral
  hyperplane class supplies a rational Lefschetz operator, and the primitive
  Hodge structures are polarizable over $\Q$.
\end{enumerate}

\begin{remark}[Polarization]
\label{rem:polarization}
Let $k\leq n$ and choose a rational K\"ahler class $[\omega]$. On the
primitive space $P^k(M,\Q)$ of \cref{rem:rational-lefschetz}, put
\[
Q_k(\alpha,\beta)
=(-1)^{k(k-1)/2}\int_M\alpha\smile\beta\smile[\omega]^{n-k}.
\]
The Hodge--Riemann bilinear relations
\cite[\S\S6.3.2 and 7.1.2]{voisin2002hodge} give, on extending
$Q_k$ complex bilinearly,
\[
Q_k(F^pP^k(M,\C),F^{k-p+1}P^k(M,\C))=0,
\qquad
\ii^{p-q}Q_k(\alpha,\overline\alpha)>0
\quad(0\ne\alpha\in P^{p,q}).
\]
Thus $Q_k$ polarizes the primitive rational Hodge structure. A real
K\"ahler class gives the corresponding real polarization on a compact
K\"ahler manifold.

For an integral $[\omega]$, the primitive lattice is the intersection of
$P^k(M,\Q)$ with $H^k(M,\Z)/\mathrm{tors}$ inside $H^k(M,\Q)$. The form
$Q_k$ takes integer values on this lattice and is nondegenerate over
$\Q$. For instance, on the quintic with $[\omega]=h$ one has
$Q_0(1,1)=5$. The homological form of the weak Lefschetz theorem and
Poincar\'e duality give $H^4(M,\Z)\cong H_2(M,\Z)\cong\Z$
\cite[Ch.~1]{voisin2003hodge2}. Let $\ell$ be the generator with
$\int_M h\smile\ell=1$. Then $h^2=5\ell$, because $\int_M h^3=5$, and the
integral Lefschetz map
\[
h\smile{-}:H^2(M,\Z)\longrightarrow H^4(M,\Z)
\]
has image of index five. The companion notes show that $\ell$ is the class
of a line on $X$ \cite[\S10.2]{hodgenotes}.

Over $\Q$, the $j$-th summand in the degree-$k$ Lefschetz decomposition
is isomorphic as a Hodge structure to $P^{k-2j}(M,\Q)(-j)$.
Here $\Q(-j)$ has type $(j,j)$, the shift in bidegree induced by $L^j$.
\end{remark}

\section{Inputs and Further Directions}
\label{sec:limits}

The analytification and GAGA theorems are inputs. The analytic proofs use the
elliptic Hodge theorem for the de Rham and Dolbeault complexes and the Dolbeault
isomorphism. The normal coordinates and the local formulas for formal adjoints
are derived in \cref{sec:kahler-identities}.

The applications use weak Lefschetz, adjunction, the Euler and normal
sequences, the Whitney formula, the cohomology of line bundles on projective
space, Poincar\'e and Serre duality, and Chern--Gauss--Bonnet. The Hopf surface
calculation also uses Hartogs' theorem and surface Riemann--Roch. The
quintic remarks use Ehresmann's theorem and the description of first-order
deformations. The
polarization statement uses the Hodge--Riemann bilinear relations. Kodaira's
embedding theorem is quoted only in \cref{rem:kodaira}.

\begin{remark}[Algebraic de Rham cohomology]
For a smooth complex algebraic variety, Grothendieck's comparison
theorem \cite{grothendieck1966derham} identifies
\[
\mathbb H^k(X,\Omega^\bullet_{X/\C})\cong H^k(\Xan,\C).
\]
For a smooth proper variety the Hodge-to-de Rham spectral sequence
degenerates at $E_1$. Deligne--Illusie prove degeneration in
characteristic zero by reduction to positive characteristic
\cite[Cor.~2.7]{deligneillusie1987}.
\end{remark}

\begin{remark}[Mixed Hodge structures]
Smoothness and properness give purity for the ordinary cohomology of a
complex algebraic variety. For smooth nonproper and for singular varieties,
ordinary cohomology generally carries a mixed Hodge structure
\cite{deligne1971hodge,deligne1974hodge}. For singular projective varieties,
intersection cohomology recovers analogues of purity and Hard Lefschetz.
\end{remark}

\begin{remark}[Cycles and variation]
For smooth projective $X$ with $M=\Xan$, the Lefschetz $(1,1)$ theorem
identifies the image of $c_1:\Pic(X)\to H^2(M,\Z)$ with
\[
\{a\in H^2(M,\Z): a_{\C}\in H^{1,1}(M)\}.
\]
The preimage formulation includes integral torsion classes. The companion
notes give the proof by the analytic exponential sequence and GAGA, and
use rational Hard Lefschetz to treat classes of curves. For the quintic
they represent the generator of $H^4(M,\Z)$ by a line
\cite[\S\S10.1--10.2]{hodgenotes}.

In a smooth projective family the Hodge filtration on the flat cohomology
bundle varies holomorphically and satisfies Griffiths transversality.
Hypersurfaces, including smooth quintics, provide examples of the
resulting period maps \cite[Ch.~6]{voisin2003hodge2}.
\end{remark}

%\clearpage
\bibliographystyle{amsalpha}
\bibliography{references}

\end{document}


%%% Local Variables:
%%% mode: LaTeX
%%% TeX-master: t
%%% End:
